\section{Witness scales and the signed estimate}
\label{sec:packing}

We finish the argument at the integer $Q$.  In this section the induction
hypothesis is needed only through \Cref{ct:separated}, as used in the
multiple-well initialization and barrier. Put $t=\log s$; a prime
below means differentiation with respect to $t$.  All measures, balls,
heights, and Gaussian kernels have the normalization of
\Cref{ex:envelope}.  In particular $I=s^2e^{-a}$, whereas the quantities
inside $\Lambda$ retain their physical units.  Write $\nu=\tau$ for the
positive projected measure in the working chart.

The purpose of \Cref{pk:witnesses,pk:reverse,pk:packing} is to turn a path of contact scales
into a small set for $\nu$.  First we find many scales with a definite
normal flux.  The height moment with exponent $p>2$ then prevents this
flux from concentrating entirely where the density of $\nu$ is large.
Finally a dyadic stopping argument makes repeated occurrences of this
property expensive in $\nu$-mass.

\subsection{Witnesses and controlled density}

\begin{lemma}[Ordered witness scales]\label{pk:witnesses}
Consider a contact path and a merger-free traversal of $[B,2B]$
provided by \Cref{ex:bands}, with $k\ge B^{-\eps_0}$ throughout the
traversal.  Given a sufficiently small fixed number $\vartheta>0$,
there are at least $cB^{1/2}$ scales $s_j$, ordered in outward time,
whose levels satisfy
\begin{equation}\label{pk:level-spacing}
 B+1<a(s_j)<2B-1,\qquad
 a(s_{j+1})-a(s_j)\ge B^{1/3}.
\end{equation}
If more than one well remains, these scales satisfy
\begin{equation}\label{pk:multi-witness}
 D_{s_j}\ge c e^{-a(s_j)},\qquad
 D_{\mathrm{tr}}(s_j)\le\vartheta e^{-a(s_j)}.
\end{equation}
If only one well remains, they satisfy
\begin{equation}\label{pk:single-witness}
 D_{s_j}\ge B^{-\eps_0}e^{-a(s_j)}.
\end{equation}
The constants are independent of the contact and of the positive
bottom scale.  The large threshold for $B$ may depend on $\eps_0$.
\end{lemma}

\begin{proof}
We first record precisely the differential information used here.
By \Cref{ex:envelope,ex:barrier}, $0\le Z\le A\le I$,
$A'-Z'\ge0$, and
\begin{equation}\label{pk:regular-path}
 |a'|+|I'|/I+|Z'|/I\le C,\qquad
 \left|(Z'/I)'\right|\le C
\end{equation}
almost everywhere.  The last estimate follows also directly by
differentiating the Gaussian source: its first two logarithmic
derivatives are bounded by Gaussian averages of fixed polynomial
weights times $\tau$, and the dilation estimate bounds these by
$CI/s^2$.  The constant background in the source is negligible since
$a\le T/2$.  At a reset $I=A$ we have $I'=A'$ almost everywhere on
the reset set.  Off that set $I'=KI$ and $a'=2-K<0$.
The weighted scalar estimate is
\begin{equation}\label{pk:weighted-scalar}
 |M|_s\le C_0(D_s+\Lambda_s)+a^{-P}I/s^2,
\end{equation}
where $P$ can be chosen as large as needed.  With several wells,
\begin{equation}\label{pk:source-transition}
 Z'+c_0\frac{(\log a)^2}{a}s^2D_{\mathrm{tr}}(s)
 \le a^{-P}I.
\end{equation}
Here $D_{\mathrm{tr}}(s)$ is the Gaussian average over the
enlarged transition set of \Cref{ex:barrier}.  This set
includes all regions where derivatives of either $\Gamma$
or $V_b$ differ from their near-well values, in particular
the cutoff annuli of $V_b$ inside the quadratic neighborhoods
of $\Gamma$; clipping tails are estimated separately.
These inequalities retain the cumulative contact condition $Z\ge0$;
no sign for $Z'$ at a later scale is asserted.

Here is a time-bin estimate that does not require monotonicity of
$a$ or of $I$.  Fix an endpoint $t_1$ of an interval on which
$a\ge A_*>2$, and put $M_* =\max_{t\le t_1}I(t)$ on that interval.
For nonnegative integers $j$ and integers $n\ge\lfloor A_*\rfloor$,
the simultaneous conditions
\[
 n\le a(t)<n+1,\qquad e^{-j-1}M_*<I(t)\le e^{-j}M_*
\]
confine $t$ to an interval of length at most one, because
\begin{equation}\label{pk:time-identity}
 t=\frac{a(t)+\log I(t)}2.
\end{equation}
The total length of the set of such times, even if it has many
components, is therefore at most one.  Summing the integral over
these sets gives
\begin{equation}\label{pk:time-bin}
 \int_{t\le t_1}a(t)^{-P}I(t)\,\dd t
 \le CM_*\sum_{j\ge0}e^{-j}
       \sum_{n\ge\lfloor A_*\rfloor}n^{-P}
 \le C A_*^{1-P}M_*.
\end{equation}
For an interval confined to $[B,2B]$, the same estimate has
$CB^{1-P}$ on the right.

Suppose first that the selected traversal has several wells.  All
earlier stages also had several wells, since the construction only
removes wells.  Its initialization gives
$Z(t_0)\le a_0^{-P}I(t_0)$.  By
\eqref{pk:source-transition} and \eqref{pk:time-bin}, the positive
variation of $Z$ up to any endpoint is at most $\epsilon M_*$,
where $\epsilon\to0$ as $a_0\to\infty$.  Since $A-Z$ is
nondecreasing, for $u\le t$ we have
\[
 A(t)\ge A(u)-Z(u)
 \ge A(u)-Z(t_0)-\int_{t_0}^t(Z')_+\,\dd v.
\]
The future envelope has the elementary property
\begin{equation}\label{pk:envelope-history}
 \max_{t_0\le u\le t}I(u)
 \le\max\left\{I(t),\max_{t_0\le v\le t}A(v)\right\}.
\end{equation}
Indeed a term realizing $I(u)$ either comes from an $A(v)$ with
$v\le t$, when its exponential weight is at most one, or comes
from a later $v$ or from the artificial endpoint, when the same
term at $t$ is larger by $e^{K(t-u)}$.  Approximating a supremum
proves the assertion if it is not attained.  Combining these facts,
and taking $a_0$ large, yields
\begin{equation}\label{pk:quasi-monotone}
 I(t)\ge \tfrac12\max_{t_0\le u\le t}I(u).
\end{equation}
This argument is unchanged at a merger: $I,A,Z$ do not jump, and
\eqref{pk:source-transition} holds on both adjacent multiple-well
stages.

Choose a traversal $[t_b,t_e]$ with $a(t_b)=B$,
$a(t_e)=2B$, and $B\le a\le2B$ between them.  Let
$M(t)=\max_{t_b\le v\le t}I(v)$ and $I_b=I(t_b)$.
On this interval
\begin{equation}\label{pk:band-positive}
 \int_{t_b}^t(Z')_+\,\dd v\le CB^{1-P}M(t),\qquad
 Z(t)\le I_b+CB^{1-P}M(t),\qquad I(t)\ge M(t)/2.
\end{equation}
Fix a small constant $\delta_1>0$ and set
$h_B=\delta_1(\log B)^2/B$.
We bound the levels whose first increasing passage has
$Z'/I<-h_B$.

Before $M$ first reaches $B^{10}I_b$, split the range of $M$ into
successive factors of two.  There are $O(\log B)$ such pieces.
Applying \eqref{pk:time-bin} on each piece shows that
\[
 \int\frac{(Z')_+}{M}\,\dd t\le CB^{1-P}\log B.
\]
Integration by parts against the continuous nondecreasing function
$M$ gives
\[
 \int\frac{-Z'}M\,\dd t
 =\frac{Z(t_b)}{M(t_b)}-\frac{Z(t)}{M(t)}
   -\int \frac{Z}{M^2}\,\dd M\le1.
\]
Consequently $\int (Z')_-/M\le2$ for large $B$.  In view of
\eqref{pk:quasi-monotone}, the time spent with $Z'/I<-h_B$ is at
most $C/h_B$.  By $|a'|\le C$, the image of these times in the
level variable has length at most
\begin{equation}\label{pk:bad-levels}
 C B/(\log B)^2=o(B).
\end{equation}
Constants here may depend on the fixed $\delta_1$.

After this growth of $M$, \eqref{pk:band-positive} implies
$Z(t)\le C(B^{-10}+B^{1-P})I(t)$.  If $Z'/I<-h_B$ at a time
with $B+1<a<2B-1$, \eqref{pk:regular-path} keeps
$Z'/I<-h_B/2$ on a following interval of length $c h_B$.
The interval stays in the traversal, and $I$ changes by at most a
fixed factor there.  It would decrease $Z$ by at least
$c h_B^2 I(t)$, exceeding the available value of $Z(t)$ for
$P>12$ and large $B$.  This contradicts $Z\ge0$.  Thus no such
bad time occurs after the stated growth.

For completeness, first increasing passages can be selected without
assuming that $a$ itself is increasing.  Use the running maximum of
$a$ starting at $t_b$.  It is Lipschitz.  Except for a null set of
levels, a first passage is a differentiability point with $a'>0$:
the image under a Lipschitz function of its nondifferentiability
set is null, and the image of its zero-derivative set is null by
the one-dimensional area formula.  First passages so obtained
are ordered.  Since off resets $a'=2-K<0$, they are resets.
The excluded levels are bounded by \eqref{pk:bad-levels}.

At every remaining passage $I=A$, so $M_s=I/s^2=e^{-a}$.
The lower bound $Z'/I\ge-h_B$ and the definition of the source
imply
\[
 C_Q\Lambda_s\le(2+h_B)e^{-a}.
\]
Apply \eqref{pk:weighted-scalar}.  Taking $C_Q$ sufficiently
large compared with $C_0$, and then $B$ large, forces
$D_s\ge c e^{-a}$.  Moreover \eqref{pk:source-transition} gives
\[
 D_{\mathrm{tr}}(s)
 \le C\left(\delta_1+\frac{B^{1-P}}{(\log B)^2}\right)e^{-a}.
\]
Choose $\delta_1$ first so that the last bound is at most
$\vartheta e^{-a}$.  This proves \eqref{pk:multi-witness} on
all but $o(B)$ of the first-passage levels in the band.

Now suppose there is one well.  If no scale in a traversal of a
level interval of length $B^{1/4}$ satisfied
\eqref{pk:single-witness}, then throughout it
$s^2D_s<B^{-\eps_0}I$.  At a reset the Gaussian horizontal
identity gives
\[
 I'=2I-s^2D_s+s^2R_s+o(I)\ge I.
\]
Also \eqref{pk:weighted-scalar}, together with $M_s=e^{-a}$,
forces $\Lambda_s\ge c e^{-a}$.  Thus $Z'\le-I$ when $C_Q$
is sufficiently large.  At a nonreset, the same scalar estimate
and $C_Q>2C_0$ instead give
\[
 Z'\le C s^2D_s+o(I)\le C B^{-\eps_0}I.
\]
These estimates explicitly include the possible geometry
contribution.  In the flat case a reset with such small tilt is
already excluded by the scalar estimate.

Put $W=Z/I$.  It is nonnegative and initially at most one.  At
resets $W'\le-1$; at nonresets $W'\le C B^{-\eps_0}$.
If the two time sets have lengths $\mathcal R,\mathcal N$,
integration gives
\[
 \mathcal R\le1+C B^{-\eps_0}\mathcal N.
\]
On the other hand $a'\le1$ at resets, while $a'=2-K$ at
nonresets.  A net increase of $B^{1/4}$ therefore requires
\[
 \mathcal R\ge B^{1/4}+(K-2)\mathcal N,
\]
a contradiction for large $B$.

In the multiple-well case select good first-passage levels
greedily with separation $2B^{1/3}$; removing an interval of
length at most $2B^{1/3}$ per choice leaves at least $cB^{2/3}$
choices.  In the one-well case use ordered, disjoint level
intervals of length $B^{1/4}$ separated by $2B^{1/3}$, and choose
a witness in each corresponding first-passage traversal.
These traversals are ordered in time.  Discarding witnesses if
necessary gives \eqref{pk:level-spacing} and the asserted number.
\end{proof}

\begin{lemma}[A controlled-density fraction]\label{pk:reverse}
Fix a finite family of nested dyadic grids with the property that
every sufficiently small ball is contained in a cube of one of
the grids with comparable diameter.  The grids may be taken
half-open, so their cubes form genuine partitions.
At each witness of \Cref{pk:witnesses} there is a cube $J$ of
one of these grids such that, with $\delta=e^{-a(s)}$,
\begin{equation}\label{pk:cube-range}
 B_s(x)\subset\operatorname{int}J,\qquad
 \ell(J)\le C s\sqrt{\log B},\qquad
 cB^{-2\eps_0}\delta\le\frac{\nu(J)}{|J|}
       \le B^{\eps_0}\delta.
\end{equation}
There are fixed constants $C_*,c_*>0$, independent of
$\eps_0$, such that
\begin{equation}\label{pk:reverse-fraction}
 \nu\left(\left\{y\in J:
    \frac{\nu(R)}{|R|}\le B^{C_*}\delta
    \text{ for every dyadic }R\subset J\text{ containing }y
    \right\}\right)
 \ge B^{-c_*\eps_0}\nu(J).
\end{equation}
Here $|J|$ denotes coordinate Lebesgue volume.  The metric volume
is uniformly comparable to it.
\end{lemma}

\begin{proof}
A finite family as stated can be constructed by using, in each
coordinate, the three dyadic grids whose shifts alternate between
the one-third positions at successive generations; this standard
Euclidean construction is described in \cite[p.~2]{HytonenKairema2012}.
At a scale
between a fixed multiple of the ball diameter and twice that
multiple, one of the three boundary sets in each coordinate is
farther than the ball diameter from the center.  Taking their
product gives the required cube.  We shall apply this property
to a larger concentric ball, so the first inclusion in
\eqref{pk:cube-range} has room to spare.

Let $V_b$ be the normal multiplier of \Cref{ex:barrier}, and put
$u=(y-x)/s$. In the scale-$s$ height coordinate $z'=z/s$, define
$v(z')=V_b(sz')/s$. Thus $D_{z'}v(z')=DV_b(sz')$, equal to the
identity when $sz'$ is near a represented well. In all integrals below,
$v$ means $v(z/s)=V_b(z)/s$, while the measures remain in the section
coordinates. Dividing the normal first-variation identity in
\Cref{ex:barrier} by $s^2$, with its Gaussian cutoff, therefore gives
\begin{equation}\label{pk:flux-lower}
 \mathcal F:=\int G_s\,v\cdot P_{vh}u\,\dd\mu
 \ge c B^{-\eps_0}\delta.
\end{equation}
For multiple wells the discarded terms are bounded by a fixed
multiple of $D_{\mathrm{tr}}(s)$, and $\vartheta$ in
\Cref{pk:witnesses} is chosen small enough to absorb them.
For one well there is no transition contribution.  Clipping and
chart tails have the arbitrarily small bounds in
\Cref{ex:barrier}, and the curved first-variation errors are
$o(B^{-\eps_0}\delta)$.  This last assertion follows from
$k\ge B^{-\eps_0}$, the moment barrier, the dilation bound,
and the physical exponential smallness in \Cref{sx:signed}.

We give the moment and tail details needed to use this flux. Fix the
witness scale $s$, and keep its wells, clipping scale, and multiplier
$V_b=V_b(s)$ fixed while the ball radius varies. Choose a fixed $p>2$
in the range \eqref{eq:lt-p-range}, and fix an exceptional proportion
small enough for the concentric dyadic comparisons below, independently
of $T$, $a$, and $R$. For $R\ge1$ in any fixed polynomial range in $a$,
the dilation bound and a Gaussian lower bound give
\[
 (Rs)^{-m}\tau(B_{cRs}(x))
 \le d_R:=C(1+R)^{K-2}\delta=o(1)
\]
on the required fixed enlargement. Its $|M|$ part identifies $Q$ by
the smooth mass average in \Cref{lt:scalar}. Choose the joint
majority-and-$L^p$ witness $\mathcal B_R$ in \Cref{lt:config}, with
that exceptional proportion. The bound $d_R=o(1)$ is uniform in
the stated range, so it lies below the corresponding fixed
small-density threshold. Its moment bound is
\[
 \left((Rs)^{-m}\int_{\pi^{-1}(B_{Rs}(x))}
       \dist(z,\mathcal B_R)^p\,\dd\mu\right)^{1/p}
 \le C(1+R)^{K/2}s\sqrt\delta.
\]
This integral includes exceptional fibers and points with zero
projection Jacobian.

Take the analogous witness $\mathcal B_1$ at scale $s$. For each of
its entries, the matched projection entry on every good fiber is
within $Cs\sqrt\delta$. The area formula and the uniform upper bound
for $J_\pi$ give these entries a fixed amount of $G_s$-weighted
$\mu$-mass. Using $|DV_b|\le C$, $|V_b|^2\le\Gamma_s$, and the
barrier $U\le s^2\delta/k$ therefore gives
\[
 \max_{\xi\in\mathcal B_1}|V_b(\xi)|^2
 \le C U+Cs^2\delta
 \le Ck^{-1}s^2\delta.
\]
This estimate is made one entry at a time, so repeated entries require
no disjoint allocation of mass. Match the centers of $\mathcal B_R$
to those of $\mathcal B_1$ through concentric dyadic dilations.
The majority sets have fixed positive relative overlaps, and the sum
of the local matching bounds is at most
$C(1+R)^{K/2}s\sqrt\delta$. For the fixed multiplier,
\[
 |v(z/s)|\le \frac C s\dist(z,\mathcal B_R)
          +\max_{\xi\in\mathcal B_R}\frac{|V_b(\xi)|}{s}.
\]
The last three bounds and the bounded normalized projected mass give
\begin{equation}\label{pk:moment-dilation}
 \left((Rs)^{-m}\int_{\pi^{-1}(B_{Rs}(x))}|v|^p\,\dd\mu
       \right)^{1/p}
 \le C(1+R)^C k^{-1/2}\sqrt\delta.
\end{equation}
All these bounds use displacements from the fixed wells. Gaussian summation of
\eqref{pk:moment-dilation} and of the energy dilation bound
gives
\begin{equation}\label{pk:flux-moments}
 \left(\int G_s|v|^p\,\dd\mu\right)^{1/p}
       \le B^{c_1\eps_0}\sqrt\delta,
 \qquad
 \int G_s(1+|u|^2)e\,\dd\mu\le C\delta.
\end{equation}
The constants $c_1$ and $C$ are fixed.  To control more distant
tails uniformly even when $s$ is very small, note that clipping
gives $|v|\le C$ and the projected mass identity and dilation
bound give $(\pi_\#\mu)_{\sqrt2s}\le C$.  Gaussian comparison
therefore gives
\[
 \int_{|u|>R}G_s|v|^p\,\dd\mu\le Ce^{-R^2/4}.
\]
At $R=C\sqrt a$, with $C$ sufficiently large, this is
smaller than $B^{-N}\delta^{p/2}$ for any prescribed fixed
$N$.  Likewise Gaussian comparison, with a slightly larger
fixed dilation to absorb $|u|^2$, bounds the energy tail by
$Ce^{-cR^2}\delta$.  These estimates cover the part beyond
the polynomial dilation range and the fixed chart cutoff;
they introduce no negative power of $s$.

Take $R_B=C_1\sqrt{\log B}$, with $C_1$ fixed sufficiently
large.  Annular Gaussian summation using
\eqref{pk:moment-dilation} makes both the flux tail outside
$B_{R_Bs}(x)$ and the tilt tail there smaller than
$o(B^{-\eps_0}\delta)$.  Choose $J$ to contain this ball in
its interior, with $\ell(J)\le C R_Bs$.
The remaining tilt has Gaussian integral at least
$cB^{-\eps_0}\delta$.  Since $G_s\le Cs^{-m}$,
\[
 \nu(J)\ge cB^{-\eps_0}\delta s^m.
\]
The upper density bound follows from the energy dilation
estimate on a ball containing $J$.  All resulting factors are
fixed powers of $\log B$, and hence are at most $B^{\eps_0}$
after increasing the threshold for $B$.  This proves
\eqref{pk:cube-range}.

Let $H$ be the union of the maximal dyadic subcubes of $J$
with density larger than $B^{C_*}\delta$.  They are disjoint,
and the upper bound in \eqref{pk:cube-range} gives
\begin{equation}\label{pk:bad-volume}
 \frac{|H|}{|J|}\le B^{\eps_0-C_*}.
\end{equation}
It is essential here to estimate mass as well as volume.
The identity $\pi_\#\mu=Q\vol_S+M$ gives
\begin{equation}\label{pk:bad-mass}
 \frac{\mu(\pi^{-1}H)}{|J|}
 \le C Q\frac{|H|}{|J|}+\frac{|M|(J)}{|J|}
 \le C B^{\eps_0-C_*}+B^{\eps_0}\delta.
\end{equation}
In particular the mass defect has not been discarded.  Writing
$\sigma=\tfrac12-\tfrac1p>0$, the Gaussian-weighted mass of
$\pi^{-1}H$ is at most
$B^{c_2\eps_0-C_*}+B^{c_2\eps_0}\delta$, after absorbing
the volume ratio $|J|/s^m$ into a power of $B^{\eps_0}$.
H\"older's inequality with exponents $p,2,1/\sigma$, together
with $|P_{vh}u|\le\sqrt e\,|u|$ and
\eqref{pk:flux-moments}, bounds the absolute flux on this set
by
\begin{equation}\label{pk:bad-flux}
 B^{c_3\eps_0}\delta
 \left(B^{c_2\eps_0-C_*}+B^{c_2\eps_0}\delta\right)^\sigma.
\end{equation}
Choose $C_*$ large, independently of sufficiently small
$\eps_0$, and then $B$ large.  Expression
\eqref{pk:bad-flux} is less than one quarter of the lower
bound in \eqref{pk:flux-lower}.  The strict inequality $p>2$
is used exactly at this step.

There remains flux at least $cB^{-\eps_0}\delta$ over
$(J\setminus H)\cap B_{R_Bs}(x)$.  Cauchy--Schwarz, the
second-moment consequence of \eqref{pk:flux-moments}, and
$|u|^2\le R_B^2$ on this core imply
\[
 cB^{-2\eps_0}\delta^2
 \le B^{c_4\eps_0}\delta\,
       C R_B^2s^{-m}\nu(J\setminus H).
\]
Thus $\nu(J\setminus H)\ge
B^{-c_5\eps_0}\delta s^m$.  Compare with
$\nu(J)\le B^{\eps_0}\delta|J|$ and absorb the remaining
fixed logarithmic factors.  This proves
\eqref{pk:reverse-fraction} with a fixed $c_*$.  The set
$J\setminus H$ is precisely the set in that formula, with
the harmless choice of strict or non-strict threshold indicated
there.
\end{proof}

\subsection{Packing and contact charging}

\begin{lemma}[Dyadic packing]\label{pk:packing}
Let $\nu$ be a finite positive measure on the union $X$ of
disjoint roots of a dyadic forest.  For a cube $R$ put
$d(R)=\nu(R)/|R|$.  Suppose that a family of marked cubes has
densities in $[\lambda_-,\lambda_+]$, where
$0<\lambda_-\le\lambda_+<\infty$, and that every marked
cube $P$ satisfies
\begin{equation}\label{pk:abstract-payment}
 \nu\{x\in P:d(R)\le Ld(P)
       \text{ for every descendant }R\ni x\}
       \ge\zeta\nu(P),
\end{equation}
where $L\ge1$ and $0<\zeta<1$.
Let $E_N$ be the union of all full cubes $F$ admitting a
marked chain
\[
 W_1\supsetneqq W_2\supsetneqq\cdots\supsetneqq W_N=F,
 \qquad d(W_{j+1})\ge S d(W_j).
\]
For $R_*>e$ and $S>LR_*$ one has
\begin{equation}\label{pk:abstract-packing}
 \nu(E_N)\le
 \left[(1-\zeta)^{N-1}
 +\frac{e}{R_*}\left(1+
      \left\lfloor\log\frac{\lambda_+}{\lambda_-}\right\rfloor
                  \right)\right]\nu(X).
\end{equation}
In particular, the full innermost witness cubes furnished by
\Cref{pk:witnesses,pk:reverse}, for a fixed band and grid, have
union $E$ satisfying
\begin{equation}\label{pk:packed-witnesses}
 \nu(E)\le C\exp(-B^{1/8})d
\end{equation}
when $\eps_0$ is sufficiently small.
\end{lemma}

\begin{proof}
We first prove the estimate for downward retracing of density.
Let
\[
 J_*=1+\left\lfloor\log(\lambda_+/\lambda_-)\right\rfloor,
 \qquad \lambda_j=\lambda_-e^j\quad(0\le j<J_*).
\]
For each $j$, stop at the first cubes, descending from the
roots, whose density is at least $\lambda_j$.  These high
stopping cubes $H$ are disjoint; a root is included if it
already qualifies.  Within each $H$ stop at the first strict
descendants $R$ with
$d(R)\le e\lambda_j/R_*$.  These low stopping cubes are
also disjoint, and
\[
 \sum_R\nu(R)
 \le \frac{e\lambda_j}{R_*}\sum_R|R|
 \le\frac{e\lambda_j}{R_*}|H|
 \le\frac e{R_*}\nu(H).
\]
Now suppose $P\supsetneqq P'$ are any two band cubes, meaning
their densities belong to $[\lambda_-,\lambda_+]$, with
$d(P')\le d(P)/R_*$.  Choose $j$ so that
$\lambda_j\le d(P)<e\lambda_j$.  The cube $P$ lies below
a high stopping cube $H$, and $P'$ lies inside a low stopping
cube below that $H$.  Hence the union $E_{\mathrm{ret}}$ of
all such low cubes has
\begin{equation}\label{pk:retrace-mass}
 \nu(E_{\mathrm{ret}})\le(eJ_*/R_*)\nu(X).
\end{equation}
This is a bound in $\nu$-mass.  Moreover, if the ancestry of
a terminal cube has such a retrace before reaching that cube,
the \emph{entire} terminal cube is in $E_{\mathrm{ret}}$.

We next construct a global forest of paying cubes.  Its first
generation consists of the first marked cubes below the roots.
For each paying cube $P$, stop at the first strict descendants
$C$ with $d(C)>Ld(P)$.  By \eqref{pk:abstract-payment},
\[
 \sum_C\nu(C)\le(1-\zeta)\nu(P).
\]
Below each such $C$ take the first marked cubes, including
$C$ itself if marked, as the next paying generation.  All cubes
in one paying generation are disjoint.  Induction gives
\begin{equation}\label{pk:survival}
 \nu(\text{union of paying generation }j)
 \le(1-\zeta)^{j-1}\nu(X).
\end{equation}

Consider a terminal $F$ with a chain $W_1,\ldots,W_N$ and
with no band retrace in its ancestry before $F$.  The first
paying cube on its ancestry contains $W_1$.  Suppose the
$j$th paying cube $P_j$ contains $W_j$.  Both cubes are band
cubes.  Absence of a retrace implies
$d(W_j)>d(P_j)/R_*$ when they differ; the lower bound needed
below also holds when they coincide.  Consequently
\[
 d(W_{j+1})\ge S d(W_j)>L d(P_j).
\]
The next density crossing, and then the next first marked cube,
therefore occur at or before $W_{j+1}$.  The $(j+1)$st paying
cube contains $W_{j+1}$.  It follows that $F$ is wholly contained
in a paying cube of generation $N$.  Thus
\[
 E_N\subset E_{\mathrm{ret}}
       \cup(\text{union of paying generation }N).
\]
Equations \eqref{pk:retrace-mass} and \eqref{pk:survival}
prove \eqref{pk:abstract-packing}.  All families are countable;
half-open cubes give the statement also for measures charging
dyadic boundaries.  Finite-depth truncation followed by
monotone convergence is an alternative justification of every
stopping construction above.

We apply \Cref{pk:packing} to all cubes which can occur in
\Cref{pk:reverse}, marking a cube if it admits at least one of
the witness labels.  A label need not be selected coherently
between different contact paths.  Each such cube satisfies
\eqref{pk:abstract-payment} with the uniform choices
\[
 \begin{aligned}
 \lambda_-&=cB^{-2\eps_0}e^{-2B},&
 \lambda_+&=B^{\eps_0}e^{-B},\\
 L&=\max\{1,c^{-1}B^{C_*+2\eps_0}\},&
 \zeta&=B^{-c_*\eps_0}.
 \end{aligned}
\]
Indeed $Ld(J)\ge B^{C_*}e^{-a}$ for any of its labels.
Take $R_* =\exp(B^{1/4})$ and
$S=\exp(B^{1/3}/2)$.  The inequality $S>LR_*$ holds for
large $B$, and $J_*=O(B)$.

There is a fixed number of grids.  A fixed fraction of the
ordered witnesses on each contact path belongs to one grid.
Retain that subsequence and reverse its order, so that cubes
are listed from outer to inner.  By \eqref{pk:regular-path}
and \eqref{pk:level-spacing}, consecutive radii differ by at
least $\exp(cB^{1/3})$.  The first two properties in
\eqref{pk:cube-range} therefore give strict nesting, with the
same contact point in every cube.  The density bounds give
the factor $S$ for consecutive densities for large $B$.
There are still $N\ge cB^{1/2}$ cubes.  Start the dyadic
forest at its maximal cubes contained in the working chart;
all witness cubes lie in this forest, and its roots are
disjoint.  Its total $\nu$-mass is at most $d$.
Thus \eqref{pk:abstract-packing} is at most
\[
 \left[\exp\{-cB^{1/2-c_*\eps_0}\}
          +CB\exp(-B^{1/4})\right]d.
\]
Choose $\eps_0$ so small that $c_*\eps_0<1/4$.
This proves \eqref{pk:packed-witnesses} after increasing the
threshold for $B$.  The estimate concerns the union of full
innermost cubes, not just the contact points or almost every
point on selected infinite chains.
\end{proof}

\begin{lemma}[Charging the contact measure]\label{pk:contact-charge}
Let $h$ and $\nu$ be finite positive measures on $\R^m$, and let
$\mathcal C$ be a set of contact centers.  For each of its
contacts suppose there are radii $0<u_x\le s_x$ with
\begin{equation}\label{pk:charge-hypotheses}
 \nu_{s_x}(x)\ge c\theta,\qquad
 \nu_{\sqrt2s_x}(x)\le C\theta,\qquad
 h_{u_x}(x)\le C\theta,
\end{equation}
where $\theta>0$ and the constants are fixed.  There is a
fixed $L_0$ such that, if $B_{L_0s_x}(x)\subset E$ for every
contact, then
\begin{equation}\label{pk:charge-bound}
 h^*(\mathcal C)\le C\nu(E).
\end{equation}
Here $E$ is Borel and $h^*$ is outer measure.  The statement
also holds for a submeasure of a contact measure whose spatial
projection is dominated by $h$.
\end{lemma}

\begin{proof}
Outside $B_{L_0s_x}(x)$ the Gaussian kernels satisfy
\[
 G_{s_x}(y-x)\le 2^{m/2}e^{-L_0^2/4}
                         G_{\sqrt2s_x}(y-x).
\]
Choose $L_0$ sufficiently large.  By
\eqref{pk:charge-hypotheses}, at least $c\theta/2$ of the
$s_x$-Gaussian integral is then inside this ball.  Since
$G_{s_x}\le Cs_x^{-m}$, the threshold balls
$B_x=B_{L_0s_x}(x)$ satisfy
\begin{equation}\label{pk:threshold-payment}
 \nu(B_x)\ge c\theta s_x^m.
\end{equation}

Apply the usual decreasing-radius selection to these balls; see
\cite[Chapter~1, Lemma~3.4]{Simon2018} for the blocking-ball property
used below.
For clarity, grouping radii in successive factors of two and
taking maximal disjoint families gives disjoint selected
balls $B_i=B_{R_i}(x_i)$ such that every original ball meets
a selected ball with radius at least half its own.  Assign
each contact to one such blocking selected ball; this gives
both
\[
 |x-x_i|<3R_i,\qquad u_x\le s_x\le2R_i/L_0.
\]
The radius condition is part of the assignment, and does not
follow merely from coverage of centers by enlarged balls.
The radii are bounded, since
$c\theta\le\nu_{s_x}(x)\le C_m s_x^{-m}\nu(\R^m)$.

Within one assigned family apply the same selection to the
derivative balls $B_{u_x}(x)$.  Obtain disjoint balls
$D_{ij}=B_{u_{ij}}(x_{ij})$ whose fivefold enlargements cover
its centers.  All the original $D_{ij}$ lie in
$B_{(3+2/L_0)R_i}(x_i)$, so comparison of ordinary volumes
gives $\sum_j u_{ij}^m\le C R_i^m$.  Positivity of the
Gaussian on each fivefold ball and
\eqref{pk:charge-hypotheses} give
\[
 h(5D_{ij})\le C\theta u_{ij}^m.
\]
Therefore the $h$ outer mass of the contacts assigned to $i$
is at most
$C\theta R_i^m\le C\nu(B_i)$, by
\eqref{pk:threshold-payment}.  The $B_i$ are disjoint and
contained in $E$, which proves \eqref{pk:charge-bound}.
This argument uses no absolute continuity of $h$, and remains
valid if the derivative radii are much smaller than the
threshold radii.
\end{proof}

\begin{proof}[Completion of the proof of \Cref{sx:signed}]
Fix $Q$ and assume the assertion for all smaller positive
integers.  Choose $C_Q$ large enough for
\Cref{ct:separated,ex:envelope,ex:barrier,pk:witnesses} and
the bottom-contact estimate in \Cref{ex:bottom}.  The
lower-integer constants are already fixed.  If $Q=1$, use
only the one-well statements; no separated estimate, function
$b$, or multiple-well construction is used.

Put $F=F_Q(T)$, $\eta=e^{-2F}$, and $a_0=F^2$, as in
\Cref{ex:envelope}.  Suppose the localized source has mass
at least $\eta d$. For each positive bottom scale use the entrance
boundary $K_\sigma$, the supported contact set $\mathcal S_\sigma$,
and the Borel threshold event $\mathcal G_\sigma$ defined in
\Cref{sec:excess}. Write $\widehat h=\widehat h_\sigma$ and
$h=h_\sigma=(\pi_X)_\#\widehat h$; both have mass $\int f$.
By \Cref{ex:bottom}, the $\widehat h$-mass of
$K_\sigma\setminus\mathcal G_\sigma$ is at most
\begin{equation}\label{pk:bottom-remainder}
 o(1)+Ce^{-a_0}d
\end{equation}
as the bottom tends to zero. We estimate the contacts in
$\mathcal G_\sigma\cap\mathcal S_\sigma$, whose $\widehat h$-mass
equals $\widehat h(\mathcal G_\sigma)$, uniformly in that bottom.

At such a contact let $\underline s$ be its derivative scale
and $s_0\ge\underline s$ a threshold equality on the
outward vertical.  In particular this includes a contact on
a vertical side below a grazing equality.  The continuous
stopping rule and localization give
\[
 \tau_{s_0}(x)=e^{-a_0},\qquad
 \tau_v(x)\le e^{-a_0}\quad(\underline s\le v\le\rho),
 \qquad h_{\underline s}(x)\le Ce^{-a_0}.
\]
The harmless cutoff errors can be absorbed in fixed constants
in these relations.  There is room for the fixed dilation
$\sqrt2s_0<\rho$: for $s\in[\rho/\sqrt2,\rho]$, the total
mass bound gives $\tau_s\le C\rho^{-m}d=o(e^{-a_0})$,
using $a_0=F_Q(T)^2=o(T)$ and
$|\log\rho|=O(F_Q(T)+\log T)$.  Thus
\eqref{pk:charge-hypotheses} holds with
$\theta=e^{-a_0}$, $u_x=\underline s$, and $s_x=s_0$.

For this periodized application, \(u_x\le s_x=s_0<\rho\), and every
fixed dilation of the balls used in the proof has a Euclidean lift
inside the fixed coordinate chart for large \(T\). The upper Gaussian
comparison holds termwise over lifts outside the torus metric ball;
at these radii \(G_s\le Cs^{-m}\), and the lower bound on fivefold
balls comes from the nearest lift. The two Euclidean selections and
their volume comparisons therefore apply in that chart, with the
radii already bounded by \(\rho\).

Each such contact has a band and a grid supplying the ordered
witnesses of \Cref{pk:witnesses,pk:reverse}. For a band $i$ and
grid $g$, let $A_{i,g}\subset\mathcal G_\sigma\cap\mathcal S_\sigma$
be the contacts admitting an eligible chain for that pair, and let
$E_{i,g}$ be the union of all its eligible full innermost cubes.
The finitely many sets $A_{i,g}$ cover
$\mathcal G_\sigma\cap\mathcal S_\sigma$; they may overlap and need
not be measurable. Each $E_{i,g}$ is Borel because the dyadic family
is countable, and \Cref{pk:packing} gives
$\nu(E_{i,g})\le C e^{-B_i^{1/8}}d$.

For a contact in $A_{i,g}$, choose any one eligible chain. Its
threshold ball lies in the innermost cube. Indeed
$a(s_0)=a_0+O(1)$, the smallest witness has level at least
$B_i\ge2a_0$, and $|a'|\le C$.
Hence
\[
 \log(s_{\min}/s_0)\ge c(B_i-a_0-O(1))\longrightarrow\infty.
\]
The cube contains $B_{s_{\min}}(x)$, so it contains
$B_{L_0s_0}(x)$ for the fixed constant in
\Cref{pk:contact-charge}.

For every subset $A\subset K_\sigma$, the induced outer measures obey
\[
 \widehat h^*(A)\le h^*(\pi_X A).
\]
Indeed, a Borel spatial set $V\supset\pi_X A$ has
$A\subset\pi_X^{-1}V$ and
$\widehat h(\pi_X^{-1}V)=h(V)$; take the infimum over such $V$.
Apply \Cref{pk:contact-charge} to the center set $\pi_X A_{i,g}$,
choosing for each center any one contact in that class and its radii.
Its threshold ball lies in $E_{i,g}$, and the Gaussian hypotheses
were verified above for the full projection $h$. Thus repeated
centers introduce no multiplicity factor, and
$h^*(\pi_X A_{i,g})\le C\nu(E_{i,g})$. The covering proof uses no
measurable choice of contacts or chains. Outer subadditivity now gives
\begin{align}
 \widehat h(\mathcal G_\sigma)
 &=\widehat h(\mathcal G_\sigma\cap\mathcal S_\sigma)\notag\\
 &\le\sum_{i,g}\widehat h^*(A_{i,g})
 \le\sum_{i,g}h^*(\pi_X A_{i,g})\notag\\
 &\le C\sum_{i,g}\nu(E_{i,g})
 \le C\sum_{i=1}^Q e^{-B_i^{1/8}}d.
 \label{pk:all-contact-bound}
\end{align}

The band choice in \Cref{ex:bands} ensures
$B_1^{1/10}\gg F^3$.  Consequently the right side of
\eqref{pk:all-contact-bound}, and also $Ce^{-a_0}d$, are
$o(\eta d)$ as $T\to\infty$.  First take $T$ large with
all fixed parameters chosen as above, and then take the
bottom small in \eqref{pk:bottom-remainder}.  These bounds
contradict $\widehat h(K_\sigma)=h(X)\ge\eta d$.

Thus the localized source has mass less than $\eta d$.
Restoring its subtractive margin, whose integral is at most
$C\eta d$, proves
\[
 \int w(L-C_Q\Lambda)\le C e^{-2F_Q(T)}d
                       \le e^{-F_Q(T)}d
\]
for sufficiently large $T$, uniformly over the stated
nonnegative unit Lipschitz weights.  This proves the
induction step and hence \Cref{sx:signed} for every positive
integer $Q$.
\end{proof}
